\chapter{Trading in a Public Cloud}\label{ch:nw:trading-in-a-public-cloud}

In one cloud provider's documentation, the \omterm{def:nw:trading-in-a-public-cloud:azid}{availability zone} called us-east-1a in one account is not necessarily the zone called us-east-1a
in another: ``AWS maps the physical Availability Zones randomly to the Availability Zone names for each AWS account''. Only the zone
identifier, a code such as use1-az1, names the same physical place in every account. Two firms that agree to meet in ``zone a'' may be
in different buildings; a trading firm that places its machines in the zone its venue calls ``a'' may be a zone boundary away from it, and
pay for it in every round trip.

Part~IV of this book is about venues that run in a \omterm{def:nw:trading-in-a-public-cloud:cloud}{public cloud}, which is where most crypto trading happens (chapter~17), and about the
cloud itself as a place to trade from. This chapter describes the cloud's geography (regions, zones and the identifiers that locate them),
the tools that place machines close together, the network and the virtualisation between a program and the wire, the noise that shared
hardware adds, and the clocks and prices. Nothing in it is measured by this book: the provider documentation is cited and dated, one
published measurement study supplies the round trips, and the rest is a labelled simulation.

\section{Regions, availability zones and zone identifiers}

\begin{definition}[Public cloud]\label{def:nw:trading-in-a-public-cloud:cloud}
A \emph{public cloud}\index{public cloud} is computing, storage and networking that a provider runs in its own data centres and rents to any
customer on demand, by the hour or the unit, through a programming interface, on hardware shared between customers unless they pay for
dedicated machines.
\end{definition}

\begin{definition}[Availability zone, availability-zone identifier]\label{def:nw:trading-in-a-public-cloud:azid}
An \emph{availability zone}\index{availability zone} is one of the isolated groups of data centres that make up a cloud region, with its own
power, cooling and networking, connected to the region's other zones by low-latency links. Its \emph{availability-zone
identifier}\index{availability-zone identifier} is the provider's code for the physical zone, the same in every account, as distinct from the
zone's name, which the provider may assign differently in each account.
\end{definition}

A cloud region (One Quant Book~3) is a geographic area; its zones are buildings or campuses within it, kilometres apart, chosen so that one
zone's failure does not take the others down. For a trading firm the consequence is chapter~9's in a new form: two machines in the same zone
talk faster than two in different zones, and the only way to know which zone a venue's servers are in is by identifier.

\begin{omfigure}
\begin{tikzpicture}[font=\footnotesize,b/.style={draw,rounded corners=2pt,minimum height=0.5cm,minimum width=1.6cm,align=center,inner sep=2pt},>=Stealth]
\node[font=\small] at (0,2.6) {account A (a firm)};
\node[font=\small] at (8.4,2.6) {account B (a venue)};
\node[font=\small] at (4.2,2.6) {physical zones};
\foreach \n/\y in {a/1.8,b/1.1,c/0.4,d/-0.3}{\node[b,fill=omDef!15] (A\n) at (0,\y) {us-east-1\n}; \node[b,fill=omDef!15] (B\n) at (8.4,\y) {us-east-1\n};}
\foreach \z/\y in {4/1.8,6/1.1,1/0.4,2/-0.3}{\node[b,fill=omThm!15] (Z\z) at (4.2,\y) {use1-az\z};}
\draw[->] (Aa) -- (Z4); \draw[->] (Ab) -- (Z6); \draw[->] (Ac) -- (Z1); \draw[->] (Ad) -- (Z2);
\draw[->] (Bb) -- (Z4); \draw[->] (Bc) -- (Z6); \draw[->] (Bd) -- (Z1); \draw[->] (Ba) -- (Z2);
\node[font=\scriptsize,align=center,omMeth] at (4.2,-1.1) {A's ``us-east-1a'' is B's ``us-east-1b'': the names disagree, the identifiers do not};
\end{tikzpicture}

\omcaption{Zone names against zone identifiers for two accounts, from the chapter's fixtures (the documented response shape of the provider's
describe-availability-zones call, with synthetic values; four of the six zones shown). The firm reads the venue's zone identifier, not its
name.}\label{fig:nw:trading-in-a-public-cloud:names}
\end{omfigure}

\begin{proposition}[Trusting names]\label{prop:nw:trading-in-a-public-cloud:names}
If a region has $n$ zones and each account's names are an independent uniform shuffle of them, two accounts' zones with the same name are the
same physical zone with probability $1/n$. If a round trip inside a zone takes $t_s$ and across zones $t_c$, choosing by name rather than by
identifier costs $(1 - 1/n)(t_c - t_s)$ per round trip on average.
\end{proposition}

\begin{proof}
Fix the venue's zone. The firm's zone with the same name is a uniformly random physical zone, so it is the venue's with probability $1/n$; the
rest of the time the round trip crosses zones and takes $t_c$ instead of $t_s$.
\end{proof}

\omcode{code/firm/cloudplan/firm_cloudplan.py}{34}{56}{Zone names resolved to identifiers from a describe-availability-zones response, and
\cref{prop:nw:trading-in-a-public-cloud:names}.}\label{lst:nw:trading-in-a-public-cloud:zones}

With the published round trips of this chapter's fourth section (median \qty{270}{\micro\second} within a zone, \qty{555}{\micro\second} across),
trusting names costs \qty{190}{\micro\second} per round trip on average in a region of three zones and \qty{238}{\micro\second} in one of six
(\cref{fig:nw:trading-in-a-public-cloud:penalty}): far more than any tuning of the machine would recover.

\begin{omfigure}
\begin{tikzpicture}
\begin{axis}[ybar,width=10cm,height=5.4cm,bar width=16pt,ymin=0,ymax=280,xtick=data,xlabel={zones in the region},
  ylabel={expected penalty (\si{\micro\second})},tick label style={font=\footnotesize},label style={font=\small},grid=major,
  grid style={gray!25},nodes near coords,nodes near coords style={font=\scriptsize},table/col sep=comma,enlarge x limits=0.2]
\addplot[fill=omDef!55,draw=omDef] table[x=n,y=penalty_us]{figdata/networks/16-trading-in-a-public-cloud/penalty.csv};
\end{axis}
\end{tikzpicture}

\omcaption{The expected round-trip penalty of choosing a zone by name instead of identifier (\cref{prop:nw:trading-in-a-public-cloud:names}),
against the number of zones in the region, with the published median round trips within and across zones. Data: \texttt{fig\_cloud.py},
\texttt{nw\_cloud.penalty\_rows()}.}\label{fig:nw:trading-in-a-public-cloud:penalty}
\end{omfigure}

The same regional logic reaches the traditional venues: CME and Google Cloud are building a private cloud region and co-location facility in
Aurora (chapter~10), so that zones and identifiers may one day matter for futures as much as they already do for crypto.

\section{Placement groups and instance selection}

\begin{definition}[Placement group]\label{def:nw:trading-in-a-public-cloud:pg}
A \emph{placement group}\index{placement group} is a customer's instruction to a cloud provider about where to put a set of its machines relative
to each other: close together for low latency (a cluster), apart for resilience (a spread), or on hardware with particular capabilities.
\end{definition}

The providers name the tool differently and promise different things. AWS's cluster \omterm{def:nw:trading-in-a-public-cloud:pg}{placement groups} keep instances in one \omterm{def:nw:trading-in-a-public-cloud:azid}{availability zone},
in the same high-bisection-bandwidth segment of its network, and recommend them for low latency; its precision time \omterm{def:nw:trading-in-a-public-cloud:pg}{placement groups} put
instances on hardware with direct access to high-precision time sources. Azure's proximity \omterm{def:nw:trading-in-a-public-cloud:pg}{placement groups} make resources ``physically located
close to each other''. Google Cloud's compact placement policies place instances close together in a zone ``on a best-effort basis'', optionally
within a maximum distance. None of them publishes a latency figure, which is why the firm measures (chapter~17).

\begin{dated}{2026-09}{Placement and time services, from provider documentation}\label{dat:nw:trading-in-a-public-cloud:docs}
\begin{itemize}
\item AWS: zone names mapped randomly per account, zone identifiers consistent across accounts; placement strategies cluster, partition, spread and
  precision time; precision time groups give an enhanced local NTP source, a PTP hardware clock for microsecond-accurate synchronisation (Linux),
  and hardware packet timestamping with nanosecond resolution, at no extra charge.
\item AWS networking: enhanced networking through SR-IOV; ENA Express raising a single flow's bandwidth from 5 to 25 Gb/s and reducing tail
  latency within a zone; \omterm{def:nw:trading-in-a-public-cloud:bare}{bare-metal instances} without virtualisation overhead.
\item Azure: proximity \omterm{def:nw:trading-in-a-public-cloud:pg}{placement groups}. Google Cloud: compact placement policies, best effort, with an optional maximum distance.
\end{itemize}
\end{dated}

Instance selection follows the same logic as chapter~8's servers: the clock and the core count of the machine, its network bandwidth, and whether
the network adapter and its offloads reach the program without a hypervisor in the way. In the cloud the choice is a line in a catalogue with an
hourly price, and the firm can try several and keep the best (chapter~17's instance lottery).

\section{Network adapters, virtualisation and bare metal}

\begin{definition}[Single-root I/O virtualisation]\label{def:nw:trading-in-a-public-cloud:sriov}
\emph{Single-root I/O virtualisation}\index{single-root I/O virtualisation} (SR-IOV) lets one physical network adapter present several virtual
functions, each assigned directly to a virtual machine, so that the machine's driver talks to the hardware without the hypervisor copying its
packets.
\end{definition}

\begin{definition}[Bare-metal instance]\label{def:nw:trading-in-a-public-cloud:bare}
A \emph{bare-metal instance}\index{bare-metal instance} is a whole physical server rented through the cloud's interface, with no hypervisor between the
customer's operating system and the hardware, and no other customer on the machine.
\end{definition}

Between a trading program on a cloud instance and the wire there are, in order: the program, the kernel or a kernel-bypass library (One Quant
Book~13), the virtual function of an SR-IOV adapter, the provider's own network card and its software, the host's physical network, and the
region's fabric. The provider documents the middle of this stack only in outline: enhanced networking through SR-IOV, ENA Express on a proprietary
transport for higher single-flow bandwidth and lower tail latency within a zone, bare metal for customers who need the hardware to themselves.
The firm can control the ends: its program and its kernel settings (chapter~3), and which instance, placement and zone it rents.

\section{Jitter and noisy neighbours}

\begin{definition}[CPU steal time, noisy neighbour]\label{def:nw:trading-in-a-public-cloud:steal}
\emph{CPU steal time}\index{CPU steal time} is the time a virtual CPU was ready to run but the hypervisor ran something else on the physical CPU; Linux
reports it as ``involuntary wait''. A \emph{noisy neighbour}\index{noisy neighbour} is another tenant of the same physical server or network whose load
takes resources (CPU, cache, memory bandwidth, network queues) from a customer's machine and so adds delay to it.
\end{definition}

The one published measurement this chapter relies on, by Hilyard and co-authors, recorded round trips between small virtual machines in one US
East region of three providers for six hours. On AWS, machines in the same subnet of one zone had a median round trip of \qty{270}{\micro\second}, a
99th percentile of \qty{395}{\micro\second} and a 99.99th of \qty{760}{\micro\second}; machines in different zones 555, 755 and
\qty{1905}{\micro\second}; and the worst round trips reached 2\,900 times the average. The model fits a lognormal to each median and 99th
percentile and adds rare tail events calibrated to the 99.99th (\cref{tab:nw:trading-in-a-public-cloud:rtt}); a third case, a shared host with a
\omterm{def:nw:trading-in-a-public-cloud:steal}{noisy neighbour}, keeps the same-zone body and adds spikes of 5 to 50 times one round trip in 500, an assumption.

\begin{table}[ht]
\centering\small
\begin{tabular}{lrrrrr}
\toprule
Case & p50 & p90 & p99 & p99.9 & p99.99 \\
\midrule
Same zone (published) & 270 & -- & 395 & -- & 760 \\
Same zone (model) & 270 & 333 & 396 & 463 & 743 \\
Cross zone (published) & 555 & -- & 755 & -- & 1\,905 \\
Cross zone (model) & 555 & 657 & 755 & 851 & 1\,908 \\
Shared host, \omterm{def:nw:trading-in-a-public-cloud:steal}{noisy neighbour} (model) & 270 & 333 & 400 & 7\,474 & 13\,940 \\
\bottomrule
\end{tabular}
\caption{Round-trip percentiles in microseconds: the published measurement (Hilyard et al., AWS, same subnet and cross AZ) and the simulation
that \texttt{firm.cloudplan} fits to it; the noisy-neighbour case is an assumption. Data: \texttt{nw\_cloud.percentiles()}.}\label{tab:nw:trading-in-a-public-cloud:rtt}
\end{table}

\begin{omfigure}
\begin{tikzpicture}
\begin{semilogyaxis}[width=11cm,height=6cm,xmin=-0.2,xmax=4.2,ymin=150,ymax=30000,xtick={0,1,2,3,4},
  xticklabels={p50,p90,p99,p99.9,p99.99},ylabel={round trip (\si{\micro\second}, log)},tick label style={font=\footnotesize},
  label style={font=\small},grid=major,grid style={gray!25},legend pos=north west,legend style={font=\footnotesize},table/col sep=comma]
\addplot[omDef,thick,mark=*,mark size=1.5pt] table[x=pos,y=same]{figdata/networks/16-trading-in-a-public-cloud/rtt.csv};
\addplot[omThm,thick,mark=square*,mark size=1.5pt] table[x=pos,y=cross]{figdata/networks/16-trading-in-a-public-cloud/rtt.csv};
\addplot[omMeth,thick,dashed,mark=triangle*,mark size=1.8pt] table[x=pos,y=noisy]{figdata/networks/16-trading-in-a-public-cloud/rtt.csv};
\legend{same zone,cross zone,{shared host, noisy neighbour}}
\end{semilogyaxis}
\end{tikzpicture}

\omcaption{Simulation fitted to a published measurement: round-trip percentiles within a zone and across zones, and on a shared host with an assumed
\omterm{def:nw:trading-in-a-public-cloud:steal}{noisy neighbour}. Crossing a zone doubles the body of the distribution; a \omterm{def:nw:trading-in-a-public-cloud:steal}{noisy neighbour} leaves the body alone and multiplies the tail. Data:
\texttt{fig\_cloud.py}.}\label{fig:nw:trading-in-a-public-cloud:rtt}
\end{omfigure}

\omcode{code/firm/cloudplan/firm_cloudplan.py}{68}{74}{The round-trip model: a lognormal fitted to a median and a 99th percentile, with rare spikes.}\label{lst:nw:trading-in-a-public-cloud:rtt}

The lesson is the one of One Quant Book~13 in a harsher setting: the median is not the problem. A strategy that races in the cloud races in the
tail, and the tail belongs to other tenants unless the firm rents the whole machine; the 99.9th percentile in the noisy-neighbour case is 16 times
the same-zone figure. Bare metal removes the neighbours on the host, not those on the network.

\section{Clocks and costs in the cloud}

Chapter~4's clocks have a cloud version. AWS documents a local time service available to every instance and, in precision time \omterm{def:nw:trading-in-a-public-cloud:pg}{placement groups},
a PTP hardware clock for ``microsecond-accurate'' synchronisation and hardware timestamps of packets with nanosecond resolution; that is enough to
timestamp the chapter's own round trips, which is what a firm must do to see its tail.

\begin{definition}[Data-transfer charge]\label{def:nw:trading-in-a-public-cloud:dtc}
A \emph{data-transfer charge}\index{data-transfer charge} is a cloud provider's fee for data moved across a boundary of its network (between zones,
between regions, or out to the internet), per gigabyte and per direction, on top of the price of the machines.
\end{definition}

\begin{dated}{2026-09}{Published cloud prices}\label{dat:nw:trading-in-a-public-cloud:prices}
AWS on-demand Linux prices in US East (N. Virginia), from the price list published on 25 September 2026: c7i.xlarge (4 vCPUs) 0.1785~USD an hour;
c7i.4xlarge (16 vCPUs, up to 12.5 Gb/s) 0.714; c6in.4xlarge (16 vCPUs, up to 50 Gb/s) 0.9072; c7i.24xlarge and the bare-metal c7i.metal-24xl (96
vCPUs, 37.5 Gb/s) 4.284 each; c7i.metal-48xl 8.568. Data within one zone is free; data crossing zones is charged 0.01~USD per gigabyte in each
direction.
\end{dated}

\begin{omfigure}
\begin{tikzpicture}
\begin{axis}[xbar stacked,width=11cm,height=4.8cm,bar width=10pt,xmin=0,xmax=7200,ytick=data,
  yticklabels from table={figdata/networks/16-trading-in-a-public-cloud/plans.csv}{plan},y dir=reverse,
  xlabel={monthly cost (USD)},tick label style={font=\footnotesize,/pgf/number format/1000 sep={\,}},label style={font=\small},grid=major,grid style={gray!25},
  legend style={font=\footnotesize,at={(0.5,-0.4)},anchor=north,draw=none},legend columns=2,table/col sep=comma]
\addplot[fill=omDef!55,draw=omDef] table[y expr=\coordindex,x=compute]{figdata/networks/16-trading-in-a-public-cloud/plans.csv};
\addplot[fill=omThm!45,draw=omThm] table[y expr=\coordindex,x=transfer]{figdata/networks/16-trading-in-a-public-cloud/plans.csv};
\legend{instances (730 hours),cross-zone data (20 TB each way)}
\end{axis}
\end{tikzpicture}

\omcaption{Three monthly plans priced from the dated catalogue of \cref{dat:nw:trading-in-a-public-cloud:prices}: four 16-vCPU instances in one zone,
the same spread over two zones exchanging 20 terabytes a month, and two bare-metal servers in one zone. Data: \texttt{fig\_cloud.py},
\texttt{nw\_cloud.plan\_costs()}.}\label{fig:nw:trading-in-a-public-cloud:plans}
\end{omfigure}

The prices are small next to chapter~9's \omterm{def:nw:colocation-products-and-how-they-are-sold:colo}{colocation}: four 16-core instances cost 2\,085~USD a month, two whole 96-core servers 6\,255~USD, where one
Nasdaq cabinet costs 7\,230~USD before its power and connections. Spreading the four instances over two zones and exchanging 20 terabytes a month adds
400~USD in \omterm{def:nw:trading-in-a-public-cloud:dtc}{data-transfer charges}, and \qty{285}{\micro\second} to every median round trip between them. In the cloud the expensive thing is not the
machine but the distance.

\begin{method}[Placing a trading system in a cloud region]\label{met:nw:trading-in-a-public-cloud:place}
\begin{enumerate}
\item Find the venue's region and, if it publishes it, the zone identifier of its servers (chapters~17 and~18); otherwise find it by measuring.
\item Resolve your own zone names to identifiers in every account you use, and place by identifier.
\item Put the latency-critical machines in one zone, in a cluster (and, if you need the clock, a precision time) \omterm{def:nw:trading-in-a-public-cloud:pg}{placement group}; consider bare
  metal for the ones that race.
\item Timestamp your own round trips with the provider's time service and measure the tail, not the median; re-measure after every maintenance
  or relaunch.
\item Price machines and data transfer together; keep cross-zone traffic for what needs resilience.
\end{enumerate}
\end{method}

\section{Tutorial: zones, tails and a bill}\label{tut:nw:trading-in-a-public-cloud:tut}

\begin{tutorial}
\textbf{Goal.} Resolve zone names, simulate round trips from a published measurement, and price three plans; no cloud account is needed.
\textbf{End state:} \cref{fig:nw:trading-in-a-public-cloud:penalty,fig:nw:trading-in-a-public-cloud:rtt,fig:nw:trading-in-a-public-cloud:plans}
and \cref{tab:nw:trading-in-a-public-cloud:rtt}.
\begin{tutsteps}
\item \textbf{Zones.} \texttt{firm\_cloudplan.load\_zones} reads the two fixtures in \texttt{data/} (the documented response shape, synthetic
  values); \texttt{same\_physical} and \texttt{locate} compare them (\cref{lst:nw:trading-in-a-public-cloud:zones}).
\item \textbf{Penalty.} \texttt{nw\_cloud.penalty\_rows()} evaluates \cref{prop:nw:trading-in-a-public-cloud:names} for 2 to 6 zones.
\item \textbf{Round trips.} \texttt{sample\_rtt} (\cref{lst:nw:trading-in-a-public-cloud:rtt}) and the three \texttt{CASES}.
\item \textbf{Bill.} \texttt{load\_catalogue} and \texttt{monthly\_cost} price \texttt{PLANS}.
\end{tutsteps}
\textbf{What to change next.} Run the chapter's code on real instances in two zones with hardware timestamps, replace the fitted case with your own
percentiles, and compare.
\end{tutorial}

\section{Build: the cloud placement plan}\label{bld:nw:trading-in-a-public-cloud:cloudplan}

\begin{build}
\bfield{Purpose} Where the firm's machines are in a cloud region, what they cost, and what their round trips look like: the base of chapters~17 to~19
and of the cloud rows of chapter~29's plan.
\bfield{Interface} \texttt{firm\_cloudplan}: \texttt{load\_zones}, \texttt{same\_physical}, \texttt{locate}, \texttt{p\_same\_name\_same\_zone},
\texttt{expected\_penalty\_us}, \texttt{RTT\_PUBLISHED}, \texttt{Rtt}, \texttt{sample\_rtt}, \texttt{Instance}, \texttt{load\_catalogue},
\texttt{monthly\_cost}.
\bfield{Rules} Place by zone identifier, never by name; prices come from the dated catalogue with their source; simulated round trips say which
measurement they are fitted to.
\bfield{Acceptance tests} \texttt{code/firm/cloudplan/tests/}: the fixtures' names and identifiers, the $1/n$ probability against a shuffle
simulation, the fit's median and 99th percentile, and a bill by hand.
\bfield{Stretch} Read live responses and price lists through the provider's interfaces; add regions and inter-region round trips.
\end{build}

\begin{omsources}
\item AWS documentation: Availability Zone IDs; placement strategies and precision time placement groups; enhanced networking, ENA Express and
Nitro bare-metal instances; Amazon VPC pricing and the public price list (September 2026).
\item Microsoft Learn, proximity placement groups; Google Cloud, placement policies.
\item O.~Hilyard, B.~Cui, M.~Webster, A.~Bangalore Muralikrishna and A.~Charapko, ``Cloudy Forecast: How Predictable is Communication Latency in the
Cloud?'', arXiv:2309.13169 (2023).
\end{omsources}

\section{Exercises}

\begin{exercise}[$\star$]\label{exo:nw:trading-in-a-public-cloud:1}
In a region of four zones, what is the probability that your ``zone b'' is the venue's ``zone b''?
\end{exercise}

\begin{exercise}[$\star$]\label{exo:nw:trading-in-a-public-cloud:2}
What does it cost a month to move 5 terabytes each way between two zones?
\end{exercise}

\begin{exercise}[$\star$]\label{exo:nw:trading-in-a-public-cloud:3}
What does \omterm{def:nw:trading-in-a-public-cloud:steal}{CPU steal time} measure, and why does a trading program care?
\end{exercise}

\begin{exercise}[$\star\star$]\label{exo:nw:trading-in-a-public-cloud:4}
Why does crossing a zone double the median round trip in the published measurement, while a \omterm{def:nw:trading-in-a-public-cloud:steal}{noisy neighbour} leaves it unchanged?
\end{exercise}

\begin{exercise}[$\star\star$]\label{exo:nw:trading-in-a-public-cloud:5}
What does a \omterm{def:nw:trading-in-a-public-cloud:bare}{bare-metal instance} remove, and what does it not?
\end{exercise}

\begin{exercise}[$\star\star$]\label{exo:nw:trading-in-a-public-cloud:6}
How would you find the zone identifier of a venue's servers if the venue does not publish it?
\end{exercise}

\begin{exercise}[$\star\star\star$]\label{exo:nw:trading-in-a-public-cloud:7}
\emph{Coding.} With \texttt{firm.cloudplan}, find the spike probability at which the noisy-neighbour case's 99.9th percentile first exceeds the cross-zone
case's.
\end{exercise}

\begin{exercise}[$\star\star\star$]\label{exo:nw:trading-in-a-public-cloud:8}
\emph{Find the flaw.} ``We launched in us-east-1a, like the venue says in its documentation, so we are in its zone.''
\end{exercise}

\section{Problem: us-east-1a Is Not Your us-east-1a}

\begin{problem}[{Weekend problem --- the price of trusting a name}]\label{pb:nw:trading-in-a-public-cloud:1}
A venue runs its matching engine in one zone of a six-zone region and tells its users the zone's name in its own account. A firm launches its machines
in the zone of the same name in its account. Use the published round trips: \qty{270}{\micro\second} within a zone, \qty{555}{\micro\second} across.

\medskip\noindent\textbf{Part I --- The names.}
\begin{enumerate}
\item With independent uniform shuffles, what is the probability that the firm is in the venue's zone?
\item What is the expected round-trip penalty, and what if the region had three zones?
\item In the chapter's fixtures, which of account A's names match account B's physically?
\item Under which name does account B see account A's us-east-1a?
\end{enumerate}

\medskip\noindent\textbf{Part II --- The tail.}
\begin{enumerate}[resume]
\item What are the modelled 99th and 99.99th percentiles within and across zones?
\item How close is the model to the published 99.99th percentiles?
\item What does a \omterm{def:nw:trading-in-a-public-cloud:steal}{noisy neighbour} do to the 99.9th percentile?
\item Which matters more for a racing strategy, the zone or the neighbour?
\end{enumerate}

\medskip\noindent\textbf{Part III --- The bill.}
\begin{enumerate}[resume]
\item What do four c7i.4xlarge instances cost a month?
\item What do 20 terabytes a month each way across zones add?
\item What do two bare-metal c7i.metal-24xl cost, and against what \omterm{def:nw:colocation-products-and-how-they-are-sold:colo}{colocation} figure should that be compared?
\item What is cheaper, a zone mistake or a bare-metal upgrade, and why is that the wrong question?
\end{enumerate}

\medskip\noindent\textbf{Part IV --- The verdict.}
\begin{enumerate}[resume]
\item State the \emph{named result}: the probability that two accounts' same-named zone is the same physical zone when names are shuffled over $n$
  zones, and the expected round-trip penalty of trusting names instead of identifiers.
\item How would the firm check its placement without the venue's help?
\item Why do providers shuffle names at all?
\item What should a venue publish, and what should a firm ask for?
\item Which parts of this problem are documentation, which measurement, which assumption?
\item How often should the placement be re-checked?
\item Does the same problem exist in a \omterm{def:nw:colocation-products-and-how-they-are-sold:colo}{colocation} hall?
\item In one sentence: what does a zone name tell you?
\end{enumerate}
\end{problem}

\section{Interview questions}

\begin{interviewq}[$\star$ \iqroles{developer}]\label{iq:nw:trading-in-a-public-cloud:1}
What is an \omterm{def:nw:trading-in-a-public-cloud:azid}{availability zone}, and why does it matter for latency?
\end{interviewq}

\begin{interviewq}[$\star\star$ \iqroles{developer}]\label{iq:nw:trading-in-a-public-cloud:2}
Your cloud trading system's median latency is fine but its 99.9th percentile is ten times worse. Where do you look?
\end{interviewq}

\begin{interviewq}[$\star\star$ \iqroles{developer}]\label{iq:nw:trading-in-a-public-cloud:3}
How do you place machines as close as possible to a venue that runs in the same cloud region?
\end{interviewq}

\begin{interviewq}[$\star\star$ \iqroles{developer, researcher}]\label{iq:nw:trading-in-a-public-cloud:4}
Can you timestamp packets accurately in a \omterm{def:nw:trading-in-a-public-cloud:cloud}{public cloud}? How?
\end{interviewq}

\begin{interviewq}[$\star\star$ \iqroles{developer}]\label{iq:nw:trading-in-a-public-cloud:5}
What does it cost to run a trading system in the cloud compared with \omterm{def:nw:colocation-products-and-how-they-are-sold:colo}{colocation}, and where do the hidden costs come from?
\end{interviewq}

\begin{interviewq}[$\star\star\star$ \iqroles{developer, trader}]\label{iq:nw:trading-in-a-public-cloud:6}
Would you run a latency-sensitive equities strategy from a \omterm{def:nw:trading-in-a-public-cloud:cloud}{public cloud}? What would have to be true?
\end{interviewq}
